% =============================================================================
% Resumo do Relatorio 2 (projeto conceitual do sistema Tag&Go).
% Base: texto canonico da biblia (secao 14.6), sem citacoes. O limiar do peso
% do varejista foi escrito como 36,6 por cento para coincidir com as Secoes 6 e 9.
% =============================================================================
\begin{center}
    \textbf{RESUMO}
\end{center}

\noindent Este relatório apresenta o projeto conceitual do Tag\&Go, sistema desenvolvido em parceria com a Riachuelo para que o cliente de moda pague pelo celular e libere sozinho a etiqueta antifurto, sem passar pelo caixa. A partir das especificações-meta revistas, agora com valor, unidade e tolerância, a análise funcional comparou duas estruturas e decompôs a função global, liberar a peça mediante pagamento, em 31 funções, seis delas básicas; entre as secundárias estão a localização do produto na loja, no nível de setor, e um ciclo logístico que devolve as etiquetas à origem para recarga e cadastro. O sistema tem dois caminhos sobre um núcleo comum: a Compra Rápida, sem instalar aplicativo, por App Clip no iPhone e página web no Android com carteira digital, e a Jornada Completa, no aplicativo da marca, com perfil, fidelidade, recomendação restrita ao estoque da loja, compra on-line e métricas de sustentabilidade. O estudo de diferenciação, a escala vertical de preços e o aproveitamento técnico alimentaram uma matriz morfológica com três alternativas; a etiqueta ativa com liberação criptográfica obteve 4,06 pontos, contra 3,85 e 3,72, resultado que se inverte quando o varejista pesa 36,6\% ou mais na decisão. A arquitetura foi descrita em cinco subsistemas e dezoito interfaces, e o desenho da etiqueta foi reformulado. Com frota de 6.854 etiquetas por loja, custo-meta de R\$~60 e meta de retorno de 99\%, a locação de R\$~2,00 por etiqueta por mês se paga se o sistema evitar 3,7 desistências de compra por dia por loja.

\vspace{1em}
\noindent\textbf{Palavras-chave:} projeto conceitual; análise funcional; etiqueta antifurto; autoatendimento; logística reversa.
